% IEEE Conference Paper - Visualization-Based Comparative Study of Classical Sorting Algorithms
\documentclass[conference]{IEEEtran}

\usepackage{graphicx}
\usepackage{amsmath}
\usepackage{amssymb}
\usepackage{booktabs}
\usepackage{array}
\usepackage{multirow}
\usepackage{microtype}
\usepackage{url}
\usepackage[noadjust]{cite}
\usepackage[hidelinks,colorlinks=false,pdfborder={0 0 0}]{hyperref}
\usepackage{float}

\hyphenation{op-tical net-works semi-conduc-tor}

\begin{document}

\title{Visualization-Based Comparative Study of Classical Sorting Algorithms}

\author{
\IEEEauthorblockN{Archit Surmal}
\IEEEauthorblockA{
Department of Computer Engineering\\
Pune Institute of Computer Technology\\
Pune, India\\
c2k231329@ms.pict.edu}
\and
\IEEEauthorblockN{Jayshree Mahajan}
\IEEEauthorblockA{
Department of Computer Engineering\\
Pune Institute of Computer Technology\\
Pune, India\\
jsmahajan@pict.edu}
\and
\IEEEauthorblockN{Pranali Navghare}
\IEEEauthorblockA{
Department of Computer Engineering\\
Pune Institute of Computer Technology\\
Pune, India\\
prnavghare@pict.edu}
}

\maketitle

% ================== ABSTRACT ==================
\begin{abstract}
A recurring gap exists between the theoretical complexity of sorting algorithms
and how they actually behave at runtime. Standard teaching approaches introduce
Big-O bounds without showing students how input structure, pivot choice, or
early termination affects real performance. This paper presents a browser-based
visualization framework that instruments eight classical sorting
algorithms---Bubble Sort, Selection Sort, Insertion Sort, Merge Sort, Quick
Sort, Heap Sort, Radix Sort, and Bucket Sort---tracking every comparison and
memory write as a typed event. The system separates algorithm execution and visual
rendering through a shared event queue. Experiments on
100-element arrays across four input distributions expose behaviors that
asymptotic notation alone cannot capture: Insertion Sort required only 108
comparisons on nearly sorted input versus 4,950 on reversed sequences---a
$45\times$ difference; Quick Sort degraded to $O(n^2)$ behavior on both nearly
sorted and reversed inputs due to its first-element pivot strategy; Merge Sort
maintained 319--542 comparisons regardless of distribution; and Radix Sort
completed with zero comparisons in every case. The tool runs entirely in a web
browser with no installation required.
\end{abstract}

\begin{IEEEkeywords}
Sorting algorithms, algorithm visualization, empirical evaluation,
event-driven architecture, educational technology, comparative analysis,
performance instrumentation
\end{IEEEkeywords}

% ================== I. INTRODUCTION ==================
\section{Introduction}

Sorting is one of the most studied problems in computer science. It serves as a
practical testbed for algorithm design ideas like divide-and-conquer, in-place
computation, and adaptive termination~\cite{cormen2009,knuth1998}. Despite
extensive theoretical treatment, students regularly struggle with the disconnect
between asymptotic bounds and actual runtime behavior. Big-O notation describes
how algorithms scale in the limit, but it hides constant factors, best-case
adaptations, and sensitivity to input distribution---all of which can dominate
performance at typical problem sizes~\cite{bentley2000,sedgewick2011}.

Existing visualization tools provide animations that help somewhat, but most
treat operation counts as secondary, focusing on visual appearance rather than
quantitative instrumentation~\cite{stasko1998,sorva2012}. Many also mix
visualization logic with algorithm code, making it harder to extend or analyze
either part independently~\cite{hundhausen2002,naps2002}.

This paper describes a framework built around three ideas: (1)~\emph{clean
separation of concerns}, where algorithms communicate only through a typed event
stream; (2)~\emph{complete instrumentation}, capturing every comparison and
memory write for offline analysis; and (3)~\emph{animated visualization},
showing the sorting process with color-coded bars at adjustable speed. Eight algorithms spanning the $O(n^2)$,
$O(n \log n)$, and $O(nk)$ complexity families are implemented, allowing direct
cross-family comparison.

% ================== II. RELATED WORK ==================
\section{Related Work}

Algorithm visualization has roots in the Balsa~\cite{brown1984} and
Tango~\cite{stasko1990} systems, which established the idea of separating
algorithm logic from its visual representation. Platforms like
VisuAlgo~\cite{halim2012} and the AlgoViz Portal~\cite{shaffer2010} cover a
broad range of algorithms but rarely expose fine-grained operation counts.
Hundhausen et al.~\cite{hundhausen2002} found that meaningful learning gains
require active student engagement rather than passive
watching~\cite{naps2002}.

For sorting in particular, Galles~\cite{galles2011} provides step-by-step
animations but minimal quantitative output. More recent visualizers by Nath et
al.~\cite{nath2021} and Patil and Marulkar~\cite{patil2021} serve educational
purposes but couple visualization tightly with algorithm code. The Sound of
Sorting~\cite{bingmann2013} adds audio feedback but prioritizes aesthetic
experience over measurement. Knuth~\cite{knuth1998}, Sedgewick~\cite{sedgewick1975},
and Bentley and McIlroy~\cite{bentley1993} all showed empirically that
real-world sorting performance is driven largely by input characteristics beyond
asymptotic class. No existing browser-based tool combines operation-level
instrumentation, a modular event-driven architecture, and quantitative-visual
analysis across multiple input distributions.

% ================== III. SORTING ALGORITHMS ==================
\section{Sorting Algorithms and Problem Statement}

The framework targets the \emph{sorting problem}: given an array $A$ of $n$
comparable elements, produce a rearrangement such that $A[1] \leq A[2] \leq
\cdots \leq A[n]$. Eight algorithms from three complexity families are
implemented and compared. Table~\ref{tab:complexity} summarizes their
theoretical properties.

\begin{table}[!t]
\caption{Theoretical Complexity of Implemented Sorting Algorithms}
\label{tab:complexity}
\centering
\setlength{\tabcolsep}{4pt}
\begin{tabular}{lcccc}
\toprule
Algorithm    & Best         & Average      & Worst        & Space      \\
\midrule
Bubble       & $O(n)$       & $O(n^2)$     & $O(n^2)$     & $O(1)$     \\
Selection    & $O(n^2)$     & $O(n^2)$     & $O(n^2)$     & $O(1)$     \\
Insertion    & $O(n)$       & $O(n^2)$     & $O(n^2)$     & $O(1)$     \\
Merge Sort   & $O(n\log n)$ & $O(n\log n)$ & $O(n\log n)$ & $O(n)$     \\
Quick Sort   & $O(n\log n)$ & $O(n\log n)$ & $O(n^2)$     & $O(\log n)$\\
Heap Sort    & $O(n\log n)$ & $O(n\log n)$ & $O(n\log n)$ & $O(1)$     \\
Radix Sort   & $O(nk)$      & $O(nk)$      & $O(nk)$      & $O(n+k)$   \\
Bucket Sort  & $O(n+k)$     & $O(n+k)$     & $O(n^2)$     & $O(n+k)$   \\
\bottomrule
\end{tabular}
\end{table}

The $O(n^2)$ family---Bubble, Selection, and Insertion Sort---works through
pairwise comparisons and in-place movement. Bubble Sort repeatedly scans the
array swapping adjacent elements that are out of order. Selection Sort finds the
minimum of the remaining unsorted section and moves it into position with one
swap per pass. Insertion Sort places each new element into its correct position
within a growing sorted prefix, which gives it $O(n)$ best-case behavior on
nearly sorted input~\cite{cormen2009}.

The $O(n \log n)$ family uses divide-and-conquer. Merge Sort splits the array
recursively and merges ordered halves, guaranteeing $O(n \log n)$ in all cases
at the cost of $O(n)$ auxiliary memory. Quick Sort partitions around a chosen
pivot; using the first element as pivot produces $O(n^2)$ behavior on sorted or
nearly sorted input~\cite{hoare1962}. Heap Sort builds a max-heap in $O(n)$
time and extracts elements one at a time, achieving $O(n \log n)$ consistently
with $O(1)$ space overhead.

The non-comparison family bypasses the $\Omega(n \log n)$ lower bound. Radix
Sort processes digits one at a time in $O(nk)$ passes. Bucket Sort distributes
elements into value-range buckets and sorts each with insertion sort, hitting
$O(n+k)$ under uniform distributions~\cite{knuth1998}.

% ================== IV. PROPOSED APPROACH ==================
\section{Proposed Approach}

The core insight behind this framework is that animation and operation counts
each reveal things the other cannot show. Watching bars move demonstrates how
data rearranges; counting comparisons and writes shows how much work the
algorithm actually does. The visualization handles the first; a separate
benchmark module handles the second.

The central design decision is an \emph{event-driven execution model}. Instead
of letting sorting algorithms write directly to the DOM or update visual
elements inline, each algorithm emits a typed event for every fundamental
operation. A shared event queue sits between event production and consumption.
The visualization renderer subscribes to this queue for animation; a benchmark
module separately instruments the same algorithm logic to collect operation
counts, with no direct coupling between the two.

Three things follow from this design. First, the event log forms a complete,
replayable execution trace, making the algorithm transparent and debuggable.
Second, adding a new algorithm requires only conforming to the event interface
--- rendering and benchmarking code stay untouched. Third, deterministic event
sequences make cross-algorithm measurement reproducible. Browser-based
delivery removes installation friction while capping practical array size at
around 500 elements for smooth animation.

% ================== V. SYSTEM ARCHITECTURE ==================
\section{System Architecture}

The framework has six loosely coupled components that communicate through the
event queue, as shown in Fig.~\ref{fig:architecture}.

\begin{figure}[!t]
\centering
\includegraphics[width=\columnwidth,keepaspectratio]{ARCHITECTURE.png}
\caption{System architecture: events flow from the algorithm engine through
the event queue to the visualization renderer and analytics module. No
component holds a direct reference to another.}
\label{fig:architecture}
\end{figure}

The \textbf{User Interface Controller} handles algorithm selection, array size
($n \in [10, 500]$), input distribution, and playback speed, and exposes pause,
step, and speed controls. The \textbf{Algorithm Engine} implements each sorting
algorithm as a JavaScript ES6 generator function that yields one typed event per
operation and pauses until the controller requests the next step. The
\textbf{Event Queue} maintains a time-ordered buffer of these
generator-produced events. The \textbf{Playback Controller} dequeues events at
the user-specified rate, notifies all subscribers, and yields to the browser
event loop to avoid blocking rendering. The \textbf{Visualization Renderer}
draws array elements as color-coded vertical bars on an HTML5 canvas: blue for
unsorted elements, red for elements currently being compared, and green for
elements in their final position. The \textbf{Analytics Module} records
comparisons, writes, and reads for each run; these counts are collected via
a browser console benchmark rather than displayed in the UI during animation.
Fig.~\ref{fig:ui} shows the complete interface.

\begin{figure}[!t]
\centering
\includegraphics[width=\columnwidth,keepaspectratio]{UI.png}
\caption{User interface during Radix Sort execution. The left panel shows
algorithm selection and theoretical complexity. The canvas renders color-coded
bars that animate the sorting process in real time.}
\label{fig:ui}
\end{figure}

Component boundaries are strictly enforced: sorting code has no references to
rendering or metrics; rendering code contains no sorting logic; analytics code
is indifferent to both. This independence was confirmed when Radix Sort was
integrated without touching any previously written component.

% ================== VI. EVENT-DRIVEN EXECUTION MODEL ==================
\section{Event-Driven Execution Model}

Five primitive event types cover all tracked operations:
\texttt{compare(i,\,j,\,result)} for element comparisons;
\texttt{swap(i,\,j)} for exchanges; \texttt{write(i,\,value)} for assignments
during merge and counting passes; \texttt{read(i)} for non-destructive
accesses; and \texttt{recursion\_start}/\texttt{end} marking recursive call
boundaries. Algorithm-specific events extend this set: \texttt{partition} for
Quick Sort, \texttt{merge} for Merge Sort, and \texttt{distribute} for Bucket
Sort.

The listing below demonstrates the generator implementation pattern:

\begin{verbatim}
function* bubbleSort(arr) {
  const n = arr.length;
  for (let i = 0; i < n - 1; i++) {
    for (let j = 0; j < n - i - 1; j++) {
      yield {type: 'compare', i: j, j: j+1};
      if (arr[j] > arr[j+1]) {
        yield { type:'swap', i: j, j: j+1 };
        [arr[j],arr[j+1]] = [arr[j+1],arr[j]];
      }
    }
  }
}
\end{verbatim}

Each \texttt{yield} pauses execution and delivers one event to the queue; the
controller steps forward by invoking \texttt{generator.next()}, ensuring a
single operation advances per cycle. All subscribers process events
synchronously, keeping the displayed animation and reported metrics in
consistent agreement.

% ================== VII. EXPERIMENTAL RESULTS ==================
\section{Experimental Setup and Empirical Results}

\subsection{Methodology}

All eight algorithms were benchmarked on arrays of 100 elements across four
input distributions, with results averaged over ten independent trials. Each
algorithm's core logic was instrumented directly, counting every comparison and
memory write without any DOM interaction. The benchmark was run on the deployed
application using the browser console, with timing measured via the
\texttt{performance.now()} API.

The four distributions are: \emph{Random} (uniformly distributed values between
5 and 100); \emph{Nearly Sorted} (a sorted sequence with 10\% random adjacent
swaps); \emph{Reversed} (strictly descending order, worst case for many
algorithms); and \emph{Duplicate-Heavy} (values drawn from a fixed set of five
values, producing approximately 70\% duplicates).

\subsection{Results and Analysis}

Table~\ref{tab:results} reports mean comparison and write counts at $n = 100$.
At this scale, JavaScript execution times were sub-millisecond across all
algorithms and are not reported; comparison and write counts are the
informative metrics here. Fig.~\ref{fig:combo} provides a visual overview.

\begin{table*}[!t]
\caption{Empirical Performance Metrics for Arrays of Size 100 Across Four
Input Distributions (Mean of 10 Trials)}
\label{tab:results}
\centering
\small
\begin{tabular}{lcccccccc}
\toprule
\multirow{2}{*}{Algorithm}
  & \multicolumn{2}{c}{Random}
  & \multicolumn{2}{c}{Nearly Sorted}
  & \multicolumn{2}{c}{Reversed}
  & \multicolumn{2}{c}{Duplicates} \\
\cmidrule(lr){2-3} \cmidrule(lr){4-5}
\cmidrule(lr){6-7} \cmidrule(lr){8-9}
& Comp & Writes & Comp & Writes & Comp & Writes & Comp & Writes \\
\midrule
Bubble     & 4950 & 4638 & 4950 &   18 & 4950 & 9890 & 4950 & 3889 \\
Selection  & 4950 &  187 & 4950 &   18 & 4950 &  104 & 4950 &  158 \\
Insertion  & 2571 & 2574 &  108 &  108 & 4950 & 5044 & 2175 & 2177 \\
Merge      &  542 &  672 &  359 &  672 &  319 &  672 &  523 &  672 \\
Quick      &  613 &  645 & 4515 &  199 & 4773 & 5010 & 1239 &  441 \\
Heap       & 1022 & 1154 & 1083 & 1272 &  944 & 1030 &  934 &  974 \\
Radix      &    0 &  200 &    0 &  200 &    0 &  300 &    0 &  200 \\
Bucket     &  160 &  100 &   85 &  100 &  215 &  100 &   95 &  100 \\
\bottomrule
\end{tabular}
\end{table*}

\textbf{$O(n^2)$ family.} Bubble Sort and Selection Sort each performed exactly
$\frac{n(n-1)}{2} = 4{,}950$ comparisons across every distribution, reflecting
their input-insensitive loop structures. Write counts diverged significantly:
Bubble Sort produced 9,890 writes on reversed input---every adjacent pair
requires a swap---but only 18 writes on nearly sorted data. Selection Sort held
writes near constant across all distributions because it performs at most one
swap per pass regardless of input order.

Insertion Sort showed the most striking adaptivity of any algorithm tested. On
nearly sorted input it needed only 108 comparisons, compared to 4,950 on
reversed data---a $45\times$ reduction. This behavior directly reflects its
inner \texttt{while} loop, which terminates as soon as it finds an element
already in the correct position. On random input it averaged 2,571 comparisons,
sitting roughly halfway between its two extremes.

\textbf{$O(n \log n)$ family.} Merge Sort maintained 319--542 comparisons
across all four distributions, with write counts fixed at exactly 672 in every
case. The consistent write count follows directly from the algorithm always
allocating and filling the same auxiliary array structure regardless of input
order.

Quick Sort showed the most notable behavior among the divide-and-conquer
algorithms. On random data it was efficient at 613 comparisons, but it degraded
sharply on both nearly sorted (4,515 comparisons) and reversed (4,773
comparisons) input. This is a direct consequence of the first-element pivot
strategy used in this implementation: when the array is already ordered or
reverse-ordered, the pivot always lands at one end of the partition, producing
$O(n^2)$ recursion. This means nearly sorted input triggers degradation just
as severely as reversed input---a behavior that asymptotic analysis does not
reveal~\cite{hoare1962}.

Heap Sort was the most stable of the three, ranging from 934 to 1,083
comparisons with no distribution causing dramatic deviation.

\textbf{Non-comparison family.} Radix Sort performed zero comparisons across
every distribution, completing each pass with digit-based bucket assignments
only. Write counts varied slightly---200 writes for two-digit values and 300
for three-digit values---reflecting the number of digit passes required. Bucket
Sort required 85--215 comparisons from the insertion sort step within each
bucket, with reversed input generating the most within-bucket work.

\begin{figure*}[!ht]
\centering
\includegraphics[width=0.9\textwidth,height=7cm,keepaspectratio]{GRAPH.png}
\caption{Comparative performance across eight sorting algorithms: comparisons
and memory writes at $n = 100$. Bubble Sort and Selection Sort exhibit the
highest comparison counts (4,950). Radix Sort achieves zero comparisons. The
$O(n \log n)$ algorithms occupy an intermediate regime.}
\label{fig:combo}
\end{figure*}

% ================== VIII. EDUCATIONAL POTENTIAL AND LIMITATIONS ==================
\section{Educational Potential and Limitations}

The framework was designed with classroom use in mind, though no formal user
study has been conducted at this stage. The animated visualization addresses
a gap that lecture-only instruction cannot easily fill: students can watch
exactly what happens when an algorithm meets a specific input distribution,
and instructors can discuss operation count differences using the benchmark
data presented in this paper.

Several design choices support learning specifically. Algorithm selection is
immediate and requires no setup, making the tool practical for in-class
demonstrations. The color coding---blue for unsorted, red for active
comparison, green for finalized---gives students a consistent visual vocabulary
across all eight algorithms. Adjustable playback speed allows an instructor to
slow down critical moments like a Quick Sort partition or a Merge Sort merge
step.

The tool is best used alongside formal instruction rather than as a replacement
for it. Watching a visualization without a prior conceptual framework tends to
produce surface-level impressions. Prediction exercises---where students
estimate operation counts before running an algorithm on a particular
distribution---are likely to produce more engagement than open-ended
exploration, based on existing research in algorithm
visualization~\cite{hundhausen2002,naps2002,stasko2009}.

Current limitations include: the browser canvas restricts smooth animation to
arrays of roughly 500 elements; JavaScript execution times at this scale are
sub-millisecond and cannot be meaningfully compared; and the framework does not
yet support side-by-side dual-algorithm comparison within a single view.

% ================== IX. FUTURE WORK ==================
\section{Future Work}

Several extensions are planned. Server-side execution in a compiled language
with event streams transmitted to the browser would enable meaningful timing
measurement at larger array sizes. The current eight-algorithm set leaves out
Shell Sort, Counting Sort, and hybrid approaches like Timsort. Enhanced
visualization modes including memory-access heat maps and simultaneous
dual-algorithm animation would allow more direct comparison. Accessibility
improvements including keyboard navigation and screen reader support would
broaden the tool's usability. A formal user study measuring learning outcomes
in a controlled classroom setting remains an important future step.

% ================== X. CONCLUSION ==================
\section{Conclusion}

This paper described a modular, event-driven browser-based framework for
visualizing and measuring eight classical sorting algorithms. The architecture
enforces strict separation among algorithm execution, visual rendering, and
analytics through a shared event queue, allowing each component to evolve
independently. Benchmark results collected directly from the implemented system
across four input distributions revealed behaviors that asymptotic notation does
not capture.

Insertion Sort dropped from 4,950 comparisons on reversed input to just 108 on
nearly sorted data---a $45\times$ difference. Quick Sort, using a first-element
pivot, degraded to near-quadratic behavior on both nearly sorted and reversed
input, showing that the worst-case trigger extends beyond just reversed
sequences. Merge Sort held its write count at exactly 672 across all
distributions. Radix Sort required zero comparisons in every case.

These findings come from real instrumented runs of the actual implementation,
making them a genuine empirical characterization of the specific algorithms as
built. The event-driven design provides a clear foundation for future work
including server-side benchmarking, larger-scale testing, and a formal
educational evaluation.

% ================== REFERENCES ==================
\begin{thebibliography}{00}

\bibitem{cormen2009}
T. H. Cormen, C. E. Leiserson, R. L. Rivest, and C. Stein,
\textit{Introduction to Algorithms}, 3rd ed.
Cambridge, MA, USA: MIT Press, 2009.

\bibitem{knuth1998}
D. E. Knuth,
\textit{The Art of Computer Programming, Volume 3: Sorting and Searching},
2nd ed. Reading, MA, USA: Addison-Wesley, 1998.

\bibitem{hundhausen2002}
C. D. Hundhausen, S. A. Douglas, and J. T. Stasko,
``A meta-study of algorithm visualization effectiveness,''
\textit{J. Visual Languages and Computing},
vol. 13, no. 3, pp. 259--290, Jun. 2002.

\bibitem{sedgewick2011}
R. Sedgewick and K. Wayne,
\textit{Algorithms}, 4th ed.
Upper Saddle River, NJ, USA: Addison-Wesley, 2011.

\bibitem{bentley2000}
J. Bentley,
\textit{Programming Pearls}, 2nd ed.
Reading, MA, USA: Addison-Wesley, 2000.

\bibitem{hoare1962}
C. A. R. Hoare,
``Quicksort,''
\textit{Computer J.},
vol. 5, no. 1, pp. 10--16, Apr. 1962.

\bibitem{stasko1998}
J. T. Stasko, J. B. Domingue, M. H. Brown, and B. A. Price,
\textit{Software Visualization: Programming as a Multimedia Experience}.
Cambridge, MA, USA: MIT Press, 1998.

\bibitem{sorva2012}
J. Sorva,
``Visual program simulation in introductory programming education,''
Ph.D. dissertation, Dept. Comput. Sci., Aalto University,
Helsinki, Finland, 2012.

\bibitem{naps2002}
T. L. Naps, G. Rossling, V. Almstrum, W. Dann, R. Fleischer,
C. Hundhausen, A. Korhonen, L. Malmi, M. McNally, S. Rodger, and
J. A. Velazquez-Iturbide,
``Exploring the role of visualization and engagement in computer science
education,''
\textit{ACM SIGCSE Bull.},
vol. 35, no. 2, pp. 131--152, Jun. 2002.

\bibitem{stasko2009}
J. Stasko and R. Kehoe,
``Using animations to learn about algorithms: An ethnographic case study,''
Georgia Institute of Technology, Atlanta, GA, USA,
Tech. Rep. GIT-GVU-09-20, 2009.

\bibitem{nath2021}
S. Nath, J. Gupta, A. Gupta, and T. Verma,
``Sorting algorithm visualizer,''
\textit{Int. Res. J. Modernization in Eng. Technol. Sci.},
vol. 3, no. 7, pp. 1--7, Jul. 2021.

\bibitem{patil2021}
P. Patil and P. S. Marulkar,
``Sorting visualizer,''
\textit{Int. J. Res. Publication and Rev.},
vol. 2, no. 6, pp. 1--6, Jun. 2021.

\bibitem{galles2011}
D. Galles,
``Data structure visualizations,''
University of San Francisco, San Francisco, CA, USA, 2011. [Online].
Available: https://www.cs.usfca.edu/\~{}galles/visualization/

\bibitem{brown1984}
M. H. Brown and R. Sedgewick,
``A system for algorithm animation,''
\textit{ACM SIGGRAPH Comput. Graph.},
vol. 18, no. 3, pp. 177--186, Jul. 1984.

\bibitem{stasko1990}
J. T. Stasko,
``Tango: A framework and system for algorithm animation,''
\textit{IEEE Computer},
vol. 23, no. 9, pp. 27--39, Sep. 1990.

\bibitem{shaffer2010}
C. A. Shaffer, M. L. Cooper, A. J. D. Alon, M. Akbar, M. Stewart,
S. Ponce, and S. H. Edwards,
``The algorithm visualization portal,''
in \textit{Proc. 41st ACM Tech. Symp. Comput. Sci. Educ.},
Milwaukee, WI, USA, Mar. 2010, pp. 440--444.

\bibitem{halim2012}
S. Halim,
``VisuAlgo: Visualising data structures and algorithms through animation,''
National University of Singapore, Singapore, 2012. [Online].
Available: https://visualgo.net/

\bibitem{bingmann2013}
T. Bingmann,
``Sound of sorting -- Audibilization and visualization of sorting
algorithms,''
Karlsruhe Institute of Technology, Karlsruhe, Germany, 2013. [Online].
Available: https://panthema.net/2013/sound-of-sorting/

\bibitem{sedgewick1975}
R. Sedgewick,
``Quicksort,''
Ph.D. dissertation, Dept. Comput. Sci., Stanford University,
Stanford, CA, USA, 1975.

\bibitem{bentley1993}
J. L. Bentley and M. D. McIlroy,
``Engineering a sort function,''
\textit{Softw. Pract. Exper.},
vol. 23, no. 11, pp. 1249--1265, Nov. 1993.

\end{thebibliography}

\end{document}
